\section{Nef plus ample is ample}

Throughout, $p:X\to\operatorname{Spec}k$ is projective, $X$ may be empty, reducible or
nonreduced, and an invertible module $N$ is \emph{nef} if its numerical class lies in the nef
cone, i.e.\ $\deg_C(N|_C)\ge0$ for every closed integral curve $C\subset X$
(Theorem~\ref{thm:native-numerical-nef-cone}). This section proves Kleiman's theorem that
nef plus ample is ample, following the proof of Lazarsfeld, \emph{Positivity in algebraic
geometry I}, Theorem~1.4.9 and Corollary~1.4.10, via the Nakai criterion. Boucksom and
Jonsson use it in this form (a nef class plus an ample class is ample) in the proof of their
lemma on relatively nef classes of test configurations.

\begin{lemma}[Positive powers of ample lines]
\label{lem:nef-ample-positive-power-ample}
\lean{EnvelopeProperty.SeshadriAmplenessCriterionNefAmple.positive_power_ample_general}
\leanok
\uses{thm:seshadri-original-strata-native-ampleness,thm:fc-original-ample-line-positive-powers,thm:seshadri-ample-top-intersection-positive}
If $A$ is ample and $a>0$, then $A^{\otimes a}$ is ample.
\end{lemma}
\begin{proof}
On the reduced closure $Z$ of each point of $X$, the restriction $A|_Z$ is ample and $Z$ is
integral, so $(A|_Z)^{\otimes a}$ is ample. Pullback commutes with tensor powers, so the
restriction of $A^{\otimes a}$ to every such $Z$ is ample, and strata ampleness
(Theorem~\ref{thm:seshadri-original-strata-native-ampleness}) concludes.
\end{proof}

\begin{lemma}[Powers of nef lines]
\label{lem:nef-ample-nef-power}
\lean{EnvelopeProperty.SeshadriAmplenessCriterionNefAmple.lineClass_tensorPower_mem_nefCone}
\leanok
\uses{thm:native-numerical-nef-cone,lem:native-numerical-line-tensor-power}
If $N$ is nef, then $N^{\otimes b}$ is nef for every $b\ge0$.
\end{lemma}
\begin{proof}
The numerical class of $N^{\otimes b}$ is $b$ times that of $N$
(Lemma~\ref{lem:native-numerical-line-tensor-power}), and the nef cone is stable under
multiplication by nonnegative reals.
\end{proof}

\begin{theorem}[Nef plus ample is ample]
\label{thm:native-nef-plus-ample}
\lean{EnvelopeProperty.SeshadriAmplenessCriterionNefAmple.isAmple_tensor_nef_of_dimension_le,
EnvelopeProperty.SeshadriAmplenessCriterionNefAmple.isAmple_tensor_nef}
\leanok
\uses{thm:seshadri-integral-ampleness-step,
thm:seshadri-original-strata-native-ampleness,
thm:seshadri-original-dimzero-ample,
thm:seshadri-original-proper-closed-map-dimension-drop,
thm:native-projective-scheme-finite-dimension-bound,
thm:seshadri-ample-top-intersection-positive,
lem:native-curve-nef-closed-restriction,
thm:native-nef-ample-mixed-euler-nonnegative,
lem:native-nef-ample-first-cut-positive,
lem:native-nef-ample-second-cut-nonnegative,
thm:nef-ample-no-positive-threshold,
thm:nef-ample-polynomial-real-positivity,
lem:nef-ample-positive-power-ample,
lem:nef-ample-nef-power,
lem:nef-ample-original-pencil-rational-scaling,
lem:nef-ample-original-pencil-derivative-nonneg,
lem:nef-ample-original-pencil-top-coefficient,
lem:nef-ample-original-pencil-threshold-identity}
\source{bj22-trivval:lem:relnef}
Let $N$ be a nef invertible module and $A$ an ample module on $X$. Then $N\otimes A$ is
ample.
\end{theorem}
\begin{proof}
We prove, by induction on $d$, the assertion for all projective $X$ with $\dim X\le d$; a
finite bound exists for every projective $X$
(Theorem~\ref{thm:native-projective-scheme-finite-dimension-bound}). For a fixed bound it
suffices to treat integral $X$: by strata ampleness
(Theorem~\ref{thm:seshadri-original-strata-native-ampleness}) we may restrict to the reduced
closure $Z$ of each point; $Z$ is integral with $\dim Z\le d$, $N|_Z$ is nef
(Lemma~\ref{lem:native-curve-nef-closed-restriction}), $A|_Z$ is ample, and
$(N\otimes A)|_Z\cong N|_Z\otimes A|_Z$.

If $d=0$, every invertible module on an integral zero-dimensional projective scheme is ample
(Theorem~\ref{thm:seshadri-original-dimzero-ample}). Assume the assertion for all projective
schemes of dimension at most $d$, and let $X$ be integral with $\dim X\le d+1$. If
$\dim X\le d$ we are done, so let $\dim X=d+1$.

For integers $a>0$ and $b\ge0$ put $B_{a,b}=N^{\otimes b}\otimes A^{\otimes a}$. For every closed
subscheme $i:Y\hookrightarrow X$ with $\dim Y\le d$, the restriction $i^*B_{a,b}$ is ample:
indeed $i^*N^{\otimes b}\cong(i^*N)^{\otimes b}$ is nef (Lemmas
\ref{lem:native-curve-nef-closed-restriction} and \ref{lem:nef-ample-nef-power}),
$(i^*A)^{\otimes a}$ is ample (Lemma~\ref{lem:nef-ample-positive-power-ample}), and the induction
hypothesis applies on $Y$.

Consider the pencil polynomials (Definition~\ref{def:nef-ample-original-pencil-polynomial})
\[
 f=P_{d+1}(\varnothing;\mathcal O_X),\qquad h=P_d(N;\mathcal O_X),\qquad
 g=P_d(A;\mathcal O_X),
\]
so that $f(t)=(N+tA)^{d+1}$, $h(t)=N\cdot(N+tA)^d$, $g(t)=A\cdot(N+tA)^d$ and $f=h+tg$
(Lemma~\ref{lem:nef-ample-original-pencil-threshold-identity}). By the scaling lemma
(Lemma~\ref{lem:nef-ample-original-pencil-rational-scaling}), for $t=a/b$ with $a,b>0$,
\[
 b^{d+1}f(t)=B_{a,b}^{d+1},\qquad b^dh(t)=N\cdot B_{a,b}^d,\qquad b^dg(t)=A\cdot B_{a,b}^d .
\]

\emph{Step 1: $f(t)>0$ at a positive rational $t=a/b$ implies that $B_{a,b}$ is ample and
$h(t)\ge0$.} Then $B_{a,b}^{d+1}>0$; every non-surjective closed integral $Y\hookrightarrow X$
has $\dim Y\le d$ (Theorem~\ref{thm:seshadri-original-proper-closed-map-dimension-drop}), so
$B_{a,b}|_Y$ is ample, and the integral step of the Nakai criterion
(Theorem~\ref{thm:seshadri-integral-ampleness-step}) shows that $B_{a,b}$ is ample. The
mixed nonnegativity theorem (Theorem~\ref{thm:native-nef-ample-mixed-euler-nonnegative})
gives $N\cdot B_{a,b}^d\ge0$, i.e.\ $h(t)\ge0$.

\emph{Step 2: $g>0$ on $(0,\infty)$.} For positive rational $t=a/b$, the first cut lemma
(Lemma~\ref{lem:native-nef-ample-first-cut-positive}), applied to $A$ and $B_{a,b}$, gives
$g(t)>0$. If $d=e+1$, the second cut lemma
(Lemma~\ref{lem:native-nef-ample-second-cut-nonnegative}) gives
$A\cdot A\cdot B_{a,b}^{e}\ge0$, which by scaling is $b^eP_e(A,A;\mathcal O_X)(t)\ge0$; hence
$g'(t)\ge0$ (Lemma~\ref{lem:nef-ample-original-pencil-derivative-nonneg}), and the same holds
trivially if $d=0$. The real-positivity theorem
(Theorem~\ref{thm:nef-ample-polynomial-real-positivity}) gives $g(t)>0$ for all real $t>0$.

\emph{Step 3: $f$ is eventually positive.} Its coefficient of $t^{d+1}$ is
$A^{d+1}>0$ (Lemma~\ref{lem:nef-ample-original-pencil-top-coefficient} and
Theorem~\ref{thm:seshadri-ample-top-intersection-positive}), so $f$ has degree $d+1$ with
positive leading coefficient.

\emph{Step 4.} By Steps 1--3 and the no-positive-threshold theorem
(Theorem~\ref{thm:nef-ample-no-positive-threshold}), $f(t)>0$ for all real $t>0$. In
particular $f(1)>0$, so $B_{1,1}\cong N\otimes A$ is ample by Step 1.
\end{proof}

\begin{corollary}[Variants]
\label{cor:native-nef-plus-ample-variants}
\lean{EnvelopeProperty.SeshadriAmplenessCriterionNefAmple.isAmple_ample_tensor_nef,
EnvelopeProperty.SeshadriAmplenessCriterionNefAmple.isAmple_tensorPower_tensor_tensorPower,
EnvelopeProperty.SeshadriAmplenessCriterionNefAmple.isAmple_tensorPower_tensor,
EnvelopeProperty.SeshadriAmplenessCriterionNefAmple.isAmple_tensor_tensorPower,
EnvelopeProperty.SeshadriAmplenessCriterionNefAmple.isAmple_pullback_tensor_nef}
\leanok
\uses{thm:native-nef-plus-ample,lem:nef-ample-positive-power-ample,lem:nef-ample-nef-power,thm:seshadri-ample-top-intersection-positive}
Let $N$ be nef and $A$ ample. Then $A\otimes N$, $N^{\otimes b}\otimes A$ and
$N^{\otimes b}\otimes A^{\otimes a}$ ($b\ge0$, $a>0$) are ample, and the restriction of
$N\otimes A$ to every closed subscheme is ample.
\end{corollary}
\begin{proof}
Apply Theorem~\ref{thm:native-nef-plus-ample} to the nef line $N^{\otimes b}$
(Lemma~\ref{lem:nef-ample-nef-power}) and the ample line $A^{\otimes a}$
(Lemma~\ref{lem:nef-ample-positive-power-ample}); use the symmetry of the tensor product, and
the fact that ampleness restricts to closed subschemes.
\end{proof}
